% !TEX TS-program = luatex
\documentclass[localFont,alternative]{documentMETADATA}

\name{Patrick}{Prell}
\tagline{Robotics \& Embedded Systems Engineer}
\socialinfo{
	\address{Sydney, NSW}
	\smartphone{+61 458 132 047}
	\email{prelly95@gmail.com}\\
	\linkedin{patprell}
	\github{Prelly95}
}

\geometry{margin=1.2cm}
\setlength{\rightcolumnlength}{15.4cm}

\begin{document}

	\makecvheader

	\makecvfooter
		{\textsc{}}
		{\textsc{Patrick Prell - Cover Letter}}
		{\thepage}

	\setlength{\parskip}{0.6em}
	\vspace{1.5em}

	% ---- Addressee and date ----
	\begin{flushleft}
		\today\\[0.8em]
		Hiring Manager\\
		Company Name\\
		Company Address
	\end{flushleft}

	\vspace{1em}

	% ---- Salutation ----
	Dear Hiring Manager,

	I am writing to express my strong interest in the Robotics Engineer position within Anduril's Maritime team. Over
	the past eight years, I have built and deployed robust embedded systems for mission-critical applications. I am
	eager to apply this experience to one of robotics' toughest challenges: ensuring autonomous surface and subsea
	vehicles operate with absolute reliability in harsh, unpredictable environments.
	My background spans complex RF hardware, safety-critical medical firmware, and field-tested UAV autonomy.

	\textbf{Signal Processing and Operational Efficiency (DroneShield):} I own the antenna calibration pipeline for an
	ultra-wideband four-antenna direction-finding system. By resolving noisy multi-channel front-end interference, I
	ensured repeatable bearing accuracy across production units. Furthermore, by coordinating across production,
	hardware, and test engineering teams, I redesigned the calibration workflow, reducing execution time from twelve
	hours to just forty minutes per unit.

	\textbf{Safety-Critical Firmware and Rigor (ResusRight):} I developed firmware for a clinical monitor compliant with
	IEC 62304 and ISO 13485 standards. Operating in an environment where failure is not an option gave me a deep,
	foundational discipline in safety architectures, field reliability, and rigorous validation qualities essential
	for autonomous maritime systems.

	\textbf{State Estimation, Computer Vision and Field Ops (CORDEL):} Working with ArduPilot and PX4 across a survey UAV
	fleet, I handled everything from airframe bring-up to field operations. I designed an Extended Kalman
	Filter (EKF) fusing IMU and RTK-GNSS data for high-precision pose and position estimation, enabling accurate
	LiDAR point-cloud generation. Additionally, I built a C++/OpenCV pipeline for optical flow pose/velocity
	estimation and spent a year leading an autonomous project in New York extracting 3D geometry from moving train
	camera feeds.

	\textbf{End-to-End Autonomy and ROS 2:} In my personal engineering work, I built a ROS 2 / PX4 SITL inspection stack for
	a VTOL platform. I use this system in the field to monitor stock and water levels on my family farm, giving me
	hands-on experience deploying ROS 2 nodes to practical, outdoor operations.

	Anduril's focus on rapid iteration and software-defined defense hardware deeply aligns with how I build systems.
	I welcome the opportunity to discuss how my state estimation, embedded firmware, and field testing experience can
	contribute to the Maritime team.

	\vspace{0.6em}

	% ---- Sign-off ----
	Kind regards,\\[0.4em]
	Patrick Prell

\end{document}
